\documentclass[11pt]{article}
\usepackage[margin=0.9in]{geometry}
\usepackage{amsmath,amssymb,amsthm,booktabs,longtable,array}
\usepackage[T1]{fontenc}
\usepackage{lmodern,microtype}
\usepackage[colorlinks=true,linkcolor=blue,urlcolor=blue]{hyperref}
\newtheorem{theorem}{Theorem}[section]
\newtheorem{lemma}[theorem]{Lemma}
\newtheorem{proposition}[theorem]{Proposition}
\newtheorem{corollary}[theorem]{Corollary}
\theoremstyle{remark}
\newtheorem{remark}[theorem]{Remark}
\newcommand{\R}{\mathbb R}
\newcommand{\C}{\mathbb C}
\newcommand{\E}{\mathbb E}
\newcommand{\Rea}{\operatorname{Re}}
\newcommand{\norm}[1]{\left\|#1\right\|}
\newcommand{\abs}[1]{\left|#1\right|}
\newcommand{\supp}{\operatorname{supp}}
\title{Toward a four-dimensional Yang--Mills construction\\
\large Notebook 03: current commutators and Gaussian cutoff clusters}
\author{Research working document}
\date{16 October 2026}
\setlength{\emergencystretch}{5em}
\begin{document}
\maketitle

\begin{abstract}
Several concrete estimates are established. First, the covariant cutoff
commutator in the deposited current argument is bounded using only
first differences, with its advertised dimension-dependent constant.
Second, sharp large-field defects of a uniformly accretive complex
Gaussian admit a convergent hard-core polymer expansion and a uniformly
small logarithm in an anchored, exponentially weighted norm. The second
proof differentiates the Gaussian covariance, never the sharp indicator.
The cluster bound also gives a direct two-source Hessian estimate and
a thermodynamic-limit criterion for the Gaussian cutoff contribution
to the quadratic coefficient. These estimates are uniform in volume
under their stated hypotheses.
The unresolved interfaces are explicit: control of the actual combined
current at complex decoupling parameters, and strict local dependence of
cutoff activities on the Yang--Mills background. No continuum construction
or mass-gap result follows without those and the other inputs listed in
Notebook 01.
\end{abstract}

\noindent\textbf{Status and scope.}
The proofs below have been checked internally, including independent
checks of the commutator algebra and the Gaussian tree summation. They
have not undergone external mathematical review. The standard forest
formula and polymer convergence theorem are cited explicitly; no novelty
is claimed for them or for these applications. References to Raw I concern
\href{https://github.com/ethanwsussman/Yang--Mills-LLM-experiment/tree/e392ef45013a581664bab76b90867170b30fe10e}
{commit \texttt{e392ef4501\allowbreak 3a581664ba\allowbreak b76b908671\allowbreak 70b30fe10e}}
of the deposited account.

\section{A cutoff commutator needing only first differences}

Fix lattice spacing $\eta>0$. For a complex covariant transport define
\[
 (T_\mu B)(x)=R_{x,\mu}B(x+\eta e_\mu),\qquad
 (T_\mu^{-1}B)(x)=R_{x-\eta e_\mu,\mu}^{-1}B(x-\eta e_\mu),
\]
\[
 \nabla_\mu=(T_\mu-I)/\eta,\qquad
 \nabla_\mu^*=(T_\mu^{-1}-I)/\eta.
\]
Here star is the bilinear transpose of the difference operator. After
complexification it need not be its Hermitian adjoint. Suppose the
transports and their inverses on the relevant stencil have norm at most
$\tau$. For a scalar cutoff $\chi$, let $s_1$ bound all its first
differences and $s_2$ all its mixed backward-forward second differences
on that stencil. Write $g$ for the supremum of the first covariant
differences of $B$, and $b$ for the supremum of $B$, there.

\begin{proposition}\label{prop:commutator}
With the exterior derivative convention
\[
 (D_W^*D_WB)_\nu
 =\sum_{\mu\ne\nu}
 (\nabla_\mu^*\nabla_\mu B_\nu-
  \nabla_\mu^*\nabla_\nu B_\mu),
\]
one has
\begin{equation}\label{eq:commutator}
 \abs{([D_W^*D_W,\chi]B)_\nu(x)}
 \le(d-1)\bigl[(2+4\tau)s_1g+(1+\tau^2)s_2b\bigr].
\end{equation}
For $1\le\tau\le5/4$, the right side is at most
$c_d(s_1g+s_2b)$, with $c_d=(2+4\tau)(d-1)$.
Thus $c_4=18$ for unitary transport and $c_4\le21$ in the indicated
complex neighborhood.
\end{proposition}

\begin{proof}
Set $a_\nu=\partial_\nu\chi$, where the scalar forward difference has
the same spacing. The exact product rules give
\[
 [\nabla_\nu,\chi]=a_\nu T_\nu,
 \qquad [\nabla_\mu^*,\chi]=-(\partial^-_\mu\chi)T_\mu^{-1},
\]
\[
 \nabla_\mu^*(a_\nu T_\nu B)
 =a_\nu\nabla_\mu^*T_\nu B
  -(\partial^-_\mu a_\nu)T_\mu^{-1}T_\nu B.
\]
The useful algebraic identity is
\begin{equation}\label{eq:no-curvature}
 \nabla_\mu^*T_\nu B
 =T_\mu^{-1}\nabla_\nu B
  -T_\mu^{-1}\nabla_\mu B-\nabla_\nu B.
\end{equation}
Multiplication by $\eta$ verifies this directly using
$T_\mu^{-1}T_\mu=I$. In particular, no two different covariant shifts
are commuted. Consequently
\begin{align*}
 [\nabla_\mu^*\nabla_\nu,\chi]B
 ={}&a_\nu\bigl(T_\mu^{-1}\nabla_\nu B
       -T_\mu^{-1}\nabla_\mu B-\nabla_\nu B\bigr)\\
 &-(\partial^-_\mu a_\nu)T_\mu^{-1}T_\nu B
  -(\partial^-_\mu\chi)T_\mu^{-1}\nabla_\nu B,
\end{align*}
which is bounded by $(1+3\tau)s_1g+\tau^2s_2b$.
For $\mu=\nu$, instead use
$\nabla_\mu^*T_\mu=-\nabla_\mu$; the bound becomes
$(1+\tau)s_1g+s_2b$. Summing the two bounds for each
$\mu\ne\nu$ proves \eqref{eq:commutator}.
All field arguments lie in the radius-$\eta$ $\ell^\infty$ stencil.
\end{proof}

For the specific operator $\Delta_{\rm cov}=\sum_\mu
\nabla_\mu^*\nabla_\mu$, the same proof gives
\begin{equation}\label{eq:laplace-comm}
 \abs{[\Delta_{\rm cov},\chi]B}
 \le d\bigl[(1+\tau)s_1g+s_2b\bigr].
\end{equation}
Application to any differently defined Laplacian must include its
additional lower-order terms.

\subsection{The scale factors in the deposited cutoff}

For $W=e^{i\eta a}U$ with $U$ unitary, the invariant Hilbert--Schmidt
norm on the Lie algebra gives
\[
 \max\{\norm{\operatorname{Ad}_W},
          \norm{\operatorname{Ad}_W^{-1}}\}
 \le e^{2\eta\norm a_{\rm op}}.
\]
Thus $\eta\sup\norm a_{\rm op}\le1/10$ implies $\tau<5/4$.
If $|a|$ denotes the deposit's normalized Hilbert--Schmidt norm for
$N\times N$ matrices, use $\norm a_{\rm op}\le\sqrt N|a|$ and
the sufficient condition $\eta\sqrt N\sup|a|\le1/10$.
The cutoff in Raw I, lines 111522--111529, is
\[
 \chi=t+s\widehat\chi_\square,\qquad
 \widehat\chi_\square(x)=\prod_\mu
 \zeta\!\left(\frac{x_\mu-y_{\square,\mu}}{M\ell_\square}\right),
 \quad 0\le\zeta\le1,\quad |\zeta'|,|\zeta''|\le5.
\]
Integral formulas for finite differences imply
$s_1\le5|s|/(M\ell_\square)$ and
$s_2\le25|s|/(M^2\ell_\square^2)$.
Set
\[
 N_y(B)=\max\left\{
 \ell_y\sup_{\widetilde\Delta(y)}|B|,
 \ell_y^2\sup_{\widetilde\Delta(y)}|\nabla_W B|\right\}.
\]
Assume the stencil is contained in $\widetilde\Delta(y)$ and the
cover has the stated comparison
$\widetilde\Delta(y)\cap\square\ne\varnothing\Rightarrow
\ell_y\le L\ell_\square$. For $M\ge2L^2$, $L\ge3$,
\[
 \frac{5\ell_y}{M\ell_\square}
 +\frac{25\ell_y^2}{M^2\ell_\square^2}
 \le\frac{55}{36}<2.
\]
The cutoff differences vanish when the enlarged stencil misses its
support. Hence
\begin{equation}\label{eq:scaled-comm}
 \ell_y^3|[D_W^*D_W,\chi]B|
 \le2c_d|s|N_y(B)
 \mathbf1_{\{\widetilde\Delta(y)\cap\square\ne\varnothing\}}.
\end{equation}
The corresponding constant from \eqref{eq:laplace-comm} is $3d$.
These constants are independent of lattice spacing, volume, and number
of scales. The geometric comparison is checked in Raw I, lines
111630--111683; it is an explicit hypothesis, not a field estimate.
Equations \eqref{eq:no-curvature}--\eqref{eq:scaled-comm} fill the local
coefficient-count step marked incomplete at lines 111615--111628.

\section{What the combined-current cancellation accomplishes}

Use the actual notation
$L_\sigma=\Delta_W-P_1(\sigma)+\mathsf K$, $Y=H(\sigma)X$, and
\[
 s_\sigma=L_\sigma A'-\mathfrak m_H(\sigma)X+F(W,\mathbb H).
\]
The definition of the cut current in Raw I, lines 111542--111543, yields
the exact identity
\begin{align}\label{eq:cut-current}
 J^\sharp[\chi]={}&\chi s_\sigma
 +[\Delta_W,\chi]A'-[P_1(\sigma),\chi]A'
 +[\mathsf K,\chi]A'\nonumber\\
 &-[D_W^*D_W,\chi]Y
 +F(W,\chi\mathbb H)-\chi F(W,\mathbb H).
\end{align}
The potentially troublesome term $\mathfrak m_H X$ cancels. Also
$Y=A'-\mathbb H$, so its weak first-order norm follows from the other
two fields. Proposition~\ref{prop:commutator} handles the displayed
second-order exterior-derivative commutator without separate Laplacian
estimates for these fields.

The projection term remains nonlocal. Its kernel is exactly
$P_1(\sigma;x,y)(\chi(y)-\chi(x))$. If the actual decoupled projection
has the required $\mathcal K_{-2}$ bound with constant $C$, its
commutator has the same bound with constant $2|s|C$ by applying the
block estimate to the two products separately. In the inherited kernel
algebra this gives
\[
 |[P_1(\sigma),\chi]A'|_{(-3),\rho}
 \le2|s|CL^{d+4}c_1|A'|_{(-1),\rho}.
\]
For the sharper spatial profile, retain $\chi(y)-\chi(x)$ instead of
its supremum. The required differentiated projection bound for general
complex $\sigma$ is an additional operator input.

\subsection{The actual carrier requirements}

The deposited transcription of CMP99b(3.37) is
$N(a)=\max\{|a|_{(-1)},|\nabla_Ua|_{(-2)}\}$; it contains no
Laplacian. The basic carrier adds a small independent current, also
without a relation $J=J(W)$. The later increment hypothesis likewise
asks only for a first-order potential increment and a combined current
increment (Raw I, lines 111779--111783 and 112378--112408).
Thus separate second-order bounds are stronger than this immediate
consumer requires.

This does not dispose of the entire domain problem. The stored domains
also impose curvature and current bounds on derived maps at every
retained triple. In particular, their fourth clause concerns
$J_{p,X}(M\!\bullet\mathbb U)$, not the independent input current.
See Raw I, lines 116925--116934 and 119480--119484. A replacement theorem
must propagate these \emph{derived combined currents} as well.

\subsection{The exact defect at complex decoupling parameters}

The earlier chart provides a useful candidate representation but must
not be identified with the later global layer by notation alone. Put
$D_0=\Delta_a-Q^*aQ$, with $a=\ell^{-2}$ the scale multiplier. For its
forced two-input solution, write
\[
 P=H_{0,\sigma}b+G_\sigma Y,
 \qquad Y=\Theta_\sigma H_{0,\sigma}b-f-V'_\sigma(P),
 \qquad Z=\Theta_\sigma P-f-V'_\sigma(P).
\]
Then $Z=(I+\Theta_\sigma G_\sigma)Y$. Define
\begin{align*}
 E_{H,\sigma}&=D_0H_{0,\sigma}
   -Q^*(q_{0,\sigma}-a)\ell^{-1},\\
 E_{G,\sigma}&=D_0G_\sigma
   -P_{0,\sigma}^{T}(I+\Theta_\sigma G_\sigma).
\end{align*}
Direct substitution proves
\begin{equation}\label{eq:operator-residual}
 D_0P=P_{0,\sigma}^{T}Z+Q^*(q_{0,\sigma}-a)\ell^{-1}b
       +E_{H,\sigma}b+E_{G,\sigma}Y.
\end{equation}
Both defects vanish at $\sigma=1$ by the undeformed inverse and block
constraint identities. The specified walk decoration at complex
$\sigma$ does not assert those identities.

More explicitly, set $X=D^J_\sigma(P)$, $H=P-H^J_\sigma X$, and
$B_\sigma X=Q^*\mathsf k_\sigma X+
\Lambda_\sigma^T\mathfrak j_J(H^J_\sigma X)$. Define
\begin{align*}
 J_{\rm sub}&=P_{0,\sigma}^{T}Z+
 Q^*(q_{0,\sigma}-a)\ell^{-1}b-B_\sigma X-\Delta'H+F(W,H),\\
 J_{\rm diff}&=(D_0-\Delta')P-
 (B_\sigma-\Delta'H^J_\sigma)X+F(W,H).
\end{align*}
Equation \eqref{eq:operator-residual} gives the precise discrepancy
\begin{equation}\label{eq:current-residual}
 J_{\rm diff}-J_{\rm sub}=E_{H,\sigma}b+E_{G,\sigma}Y.
\end{equation}
Matching the global operators would add the differences between
$L_\sigma$ and $D_0-\Delta'$, and between $\mathfrak m_H(\sigma)$ and
$B_\sigma-\Delta'H^J_\sigma$. These matches are not proved here.
An ordinary weak walk estimate does not bound the differentiated
defects uniformly as $\eta\downarrow0$. A local parametrix identity or
differentiated walk estimate is the remaining concrete task.

\subsection{A possible replacement for the derived-current bound}

There is a separate first-derivative route for a configuration which
is actually critical under a constraint. It must not be applied to
an arbitrary independent-current pair.

\begin{proposition}\label{prop:derived-current}
Let $Q_s$ be the linearized constraint at a configuration $U$, with
right inverse $E_s$ on its actual constraint space, $Q_sE_s=I$.
Use the bilinear pairings, lattice measures, and transposes of that
constraint. Suppose the physical current has the Wilson form
$J(U)=D_U^TF_U$, and constrained criticality gives
$J(U)=Q_s^T\lambda$. Then
\begin{equation}\label{eq:derived-current}
 J(U)=\mathcal R_s(U)F_U,\qquad
 \mathcal R_s(U)=Q_s^T(D_UE_s)^T.
\end{equation}
If instead $J(U)=Q_s^T\lambda+r$, the exact additional term is
$(I-Q_s^TE_s^T)r$.
\end{proposition}

\begin{proof}
Multiply the criticality equation by $E_s^T$ and use
$E_s^TQ_s^T=I$. Substitute $J=D_U^TF_U$, then multiply by $Q_s^T$.
With a residual, the same substitution gives the stated correction.
\end{proof}

For the earlier chart, $QH_0=\ell^{-1}$, so the correctly normalized
candidate is $E_s=H_{0,s}\ell_s$, provided this is the actual composite
constraint. This identity uses a first derivative of $E_s$.
For example, suppose blocks in the retained structure have side $h_y$
and metric $d_s$, and for $f$ supported in block $z$,
\[
 \sup_{\Delta_s(y)}|\mathcal R_sf|
 \le C_R h_y^{-1}e^{-\delta d_s(y,z)}\sup|f|.
\]
If $|F_U|_{\Delta_s(z)}\le a_z$ and
$a_z\le C_a a_y e^{\kappa_s d_s(y,z)}$, with
$0\le\kappa_s<\delta$, then
\[
 |J(U)|_{\Delta_s(y)}
 \le C_RC_a h_y^{-1}a_y
       \sum_z e^{-(\delta-\kappa_s)d_s(y,z)}.
\]
Thus a uniform summability bound gives the intended curvature divided
by block-side estimate. For the deposited fourth clause, its two
scales require $h_y=\ell_{p\wedge\lambda(y)}$ and a curvature
profile proportional to
$\omega(y)\varepsilon_{m(y)}\ell_{m(y)}^{-2}$.
The profile must be controlled in the \emph{retained structure's}
metric, not merely in the original hierarchy's metric.

Raw I, lines 128807--128829 and 128923--128930, treats the needed
multilevel current estimate as an input. Equation
\eqref{eq:derived-current} supplies a possible route to that input:
identify each composite right inverse and bound its differentiated
transpose kernel and profile convolution uniformly. For increments
on a fixed structure, one also needs
\[
 \delta J=\mathcal R_s(U_1)\delta F+
          [\mathcal R_s(U_1)-\mathcal R_s(U_0)]F_{U_0}.
\]
A holomorphic kernel bound on a larger ball would control the second
term by Cauchy's formula. Neither that domain nor the uniform
composite-kernel statement is proved by the algebraic identity alone.

\subsection{Differentiating the first piece of a walk}

There is a concrete way to estimate the defects in
\eqref{eq:operator-residual} without assuming they vanish. In the
earlier chart put
\[
 \mathcal A=\Delta+DD^*,\quad P_1=DPD^*,\quad
 D_0=\mathcal A-P_1.
\]
For a local inverse $G_\square$ of
$\mathcal A-P_{1,\square}+B_\square$, with
$B_\square=Q_\square^*a_\square Q_\square$, let
$S_\square=h_\square G_\square h_\square$.
Assume the local and global differential coefficients agree on the
cutoff's required stencil. The local inverse equation gives exactly
\begin{align}\label{eq:head}
 D_0S_\square={}&h_\square^2
 +[\mathcal A,h_\square]G_\square h_\square
 -P_1S_\square\nonumber\\
 &+h_\square P_{1,\square}G_\square h_\square
 -h_\square B_\square G_\square h_\square.
\end{align}
This follows simply by moving $\mathcal A$ through $h_\square$.
The global $Q^*aQ$ term has already canceled in $D_0$; it is absent
from this formula.

Both exterior-derivative parts of the commutator can be handled with
first differences. In addition to Proposition~\ref{prop:commutator},
the scalar-to-bond convention $(DD^*B)_\nu=
\sum_\mu\nabla_\nu\nabla_\mu^*B_\mu$ gives
\begin{align*}
 |([DD^*,\chi]B)_\nu|
 \le{}&[(1+\tau)+(d-1)(1+\tau+2\tau^2)]s_1g\\
 &+[1+(d-1)\tau^2]s_2b.
\end{align*}
For mixed indices, use
$\nabla_\nu T_\mu^{-1}=T_\nu\nabla_\mu^*
 +\nabla_\nu-\nabla_\mu^*$ and
$\nabla_\mu^*=-T_\mu^{-1}\nabla_\mu$ in the product rule.
The diagonal term has the bound already proved. Since
$\mathcal A=D^*D+DD^*+\Delta'$, the remaining local term to bound
is $[\Delta',h_\square]G_\square h_\square$.
This does not identify the separate global symbol $\Delta_W$ with
the component Laplacian of Section 1.

For clarity, the summation principle is elementary. Suppose a
decorated inverse is an absolutely convergent sum of products
$m_\sigma(\gamma)S_{\gamma_0}R_{\gamma_1}\cdots R_{\gamma_n}$.
If, in the intended weighted block norms, the sums of the
$D_0S_\square$ majorants from the five terms of \eqref{eq:head}
and the remainder majorants are bounded by $H$ and $r$,
and $|m_\sigma(\gamma)|\le u^{n+1}$ with $ur<1$, then
\[
 \norm{D_0G_\sigma}\le\frac{uH}{1-ur}.
\]
Apply $D_0$ to the first piece, use \eqref{eq:head}, and sum the
geometric series. This requires a bound on those five head terms
in the stated norms; it does not infer one from weak convergence.

The relevant head is present in the deposited parametrix at Raw I,
lines 25577--25610. After expansion, $H_0$ and $\widetilde G^J$
start with this inverse. Their current insertion $\Theta_J$ maps the
weak field space boundedly into the forcing/current space. The analogous
$G_\pi^J$ insertion contains $\psi+D\varphi$, so differentiating its
head can expose the extra derivative $D\varphi$. Consequently this
repair does not automatically establish the previously withheld
Laplacian bound for $H^J$. Its useful target is specifically the
defects $E_H,E_G$, subject to uniform head bounds and the separate
identification with the selected global layer.

\subsection{The averaging head: measures and a proved flat bound}

The last term of \eqref{eq:head} needs an actual averaging estimate.
Its scale cannot be read from the printed diagonal coefficient while
silently changing the transpose convention. Let $F=\eta^d I$ be the
fine-lattice measure and $M_c=\operatorname{diag}(\ell_c^d)$ the
coarse-volume measure. For the discrete matrix $Q$, set
\[
 Q^T_{\rm count}=F^{-1}Q^{\mathsf t},\qquad
 Q^T_{\rm vol}=F^{-1}Q^{\mathsf t}M_c.
\]
The transpose is bilinear, including the Lie-algebra fiber. Exactly,
\begin{equation}\label{eq:measure-conversion}
 Q^T_{\rm count}\operatorname{diag}(a_n\ell_n^{d-2})Q
 =Q^T_{\rm vol}\operatorname{diag}(a_n\ell_n^{-2})Q.
\end{equation}
Raw I, lines 12756--12760, specifies the volume pairing, whereas
lines 23961--23968 print the first diagonal coefficient in spaces
denoted $\ell^2$ without specifying their measures there. Thus
\eqref{eq:measure-conversion} is a necessary convention dictionary,
not a claim of an inconsistency in the primary paper. If the fine
pairing also changes to counting measure, its factor must be converted
as well. The dimensionless sequence $a_n$ must retain its own bound.

\begin{proposition}[Flat bond averaging]\label{prop:flat-average}
In dimensionless fine coordinates let $M=L^n$ and
\[
 (Q_M A)_\mu(c)=M^{-(d+1)}
       \sum_{x\in B_M(c)}\sum_{t=0}^{M-1}A_\mu(x+t e_\mu).
\]
Writing $q_M(c,b)$ for its scalar matrix coefficients, one has
\[
 q_M(c,b)\ge0,\quad \sum_bq_M(c,b)=1,\quad
 q_M(c,b)\le M^{-d},\quad \sum_cM^dq_M(c,b)=1.
\]
Consequently, at physical coarse spacing $\ell=M\eta$,
\[
 \norm{Q_M^T{}_{\rm vol}\,a\ell^{-2}Q_M}_{\infty\to\infty}
 \le |a|\ell^{-2}.
\]
Restricting the coarse rows or taking a Dirichlet restriction preserves
these upper bounds.
\end{proposition}

\begin{proof}
The formula at one step is the flat main term of Raw I,
lines 24355--24369; its correction brackets vanish at the identity.
Iteration gives the displayed expression: successive block and line
coordinates supply independently the base-$L$ digits of $x$ and $t$.
The row sum counts $M^{d+1}$ pairs. For a fixed fine bond, the
transverse coordinates select one coarse block. The longitudinal
coefficient is $M^{-(d+1)}k_M(s)$, where $k_M(s)$ counts pairs
$(x_\mu,t)$ in an $M$-block and $\{0,\ldots,M-1\}$ with sum $s$.
Every residue modulo $M$ has exactly $M$ such pairs, so
$\sum_zk_M(s-Mz)=M$. This proves the weighted column identity.
The final estimate follows by the two positive sums; deleting rows
or columns only decreases them.
\end{proof}

Here is a mixed-scale consequence with all needed geometry stated.
Suppose at most $N_{\rm lev}$ coarse levels contribute to an output
of $hB_\square$, the input and output scales on each averaging support
are within a factor $R$ of its coarse scale $\ell_c$, and
$\sup_n|a_n|\le a_*$. In the norms
$|A|_{(-1)}=\sup_b\ell_b|A_b|$ and
$|f|_{(-3)}=\sup_b\ell_b^3|f_b|$,
\begin{equation}\label{eq:flat-mixed-average}
 |hB_\square A|_{(-3)}
 \le\norm h_\infty N_{\rm lev}a_*R^4|A|_{(-1)}.
\end{equation}
Indeed $|QA(c)|\le R\ell_c^{-1}|A|_{(-1)}$, while
$\ell_b^3\ell_c^{-3}\le R^3$ and the volume-transpose
coefficients sum to at most one at each level. The indicated local
geometry in Raw I, lines 25742--25747, has only levels $j,j+1$;
with the support containment at lines 25479--25483 one may use
$N_{\rm lev}=2$, $R=L$ there. If the local inverse bound is $C_G$,
this flat head term is at most
$2a_*L^4 C_G\norm h_\infty^2$. Conversion to the global hierarchy
still requires comparable local and global weights at outer seams.
An exponential weight costs at most $e^{\kappa D_Q}$ when each
averaging support has diameter $D_Q$ in the retained metric.

For a complex background, a sufficient replacement is the genuinely
normalized multilevel kernel estimate
\begin{equation}\label{eq:complex-average-needed}
 \sum_b\norm{q_n(c,b)}\le C_{\rm row},\qquad
 \sup_b\sum_{c\text{ at level }n}
    (\ell_n/\eta)^d\norm{q_n(c,b)}\le C_{\rm col},
\end{equation}
uniformly in $n$ and on the selected complex domain.
Then \eqref{eq:flat-mixed-average} gains only the factor
$C_{\rm row}C_{\rm col}$. Raw I, lines 46981--46984, explicitly
includes $L^{-d}$ in the one-step kernel formula; its later printed
bound $|Q_j(U;c,b)|\le4$ must therefore not be read as an unweighted
matrix bound. The multilevel normalization, support incidence, and
complex domain of that source estimate remain to be matched here.
A bound obtained by multiplying a fixed constant at every scale
would not suffice.

The support count in this proposed completion is elementary. A
coarse bond's two endpoint blocks contain $2M^d$ fine sites and at
most $2dM^d$ relevant positive fine bonds. A fixed fine bond can be
read by at most $2d$ coarse bonds at that level. Consequently a
normalized bound $q_n=M^{-d}K_n$, $\norm{K_n}\le4$, would give
$C_{\rm row},C_{\rm col}\le8d$, without multiplying that constant
over the hierarchy. This is conditional on the normalized bound on
the selected complex domain, not its verification.

There is a second, restricted-patch route to the same primitive.
For an operator from spacing $s$ to $t$, define
\[
 \norm P_{\rm av}=\max\left\{
 \sup_c\sum_b\norm{P_{cb}},\quad
 \sup_b\sum_c(t/s)^d\norm{P_{cb}}\right\}.
\]
This norm is submultiplicative across successive scales because the
volume ratios multiply. The corrected one-step operator has
$\norm{Q(1)}_{\rm av}=1$ by Proposition~\ref{prop:flat-average}.
Fixed-dimensional holomorphy on each double block, with a uniform
coefficient bound on a fixed small ball, gives by Cauchy's formula
\[
 \norm{Q(e^{iB})}_{\rm av}\le1+C_L\norm B_\infty.
\]
If the intermediate backgrounds on the entire read patch satisfy
$\norm{\log\mathcal M^p(U)}_\infty\le C_{\rm av}bL^{p-n}$
for $0\le p<n$, within that ball, composition proves
\[
 \norm{Q_n(U)}_{\rm av}
 \le\prod_{p=0}^{n-1}(1+C_LC_{\rm av}bL^{p-n})
 \le\exp\!\left(\frac{C_LC_{\rm av}b}{L-1}\right).
\]
Every individual coefficient then obeys the normalized bound
$\norm{q_n(c,b)}\le L^{-nd}\norm{Q_n(U)}_{\rm av}$.
The deposited use of CMP98 Proposition 6 at Raw I, lines
17806--17810, has precisely the amplitude-times-block-length form
needed for this geometric sequence. Its full primary domain and
locality hypotheses have not been verified here. Applying it to
the actual head still needs one common regular gauge on the whole
read patch, the retained lower bound on local scales, and the
local/global weight comparison. The flat norm one belongs to the
corrected operator of lines 24355--24365, not to the separate rooted
$Q^\flat=Q_{\rm st}-D_1aT$ used in Notebook 05.

\subsection{The local curvature correction in the head}

The stated weak local estimate in Raw I, lines 13617--13620, is
\[
 |\Delta'A|_{(-3)}\le C_\Delta(a_0+a_1)|A|_{(-1)}.
\]
It implies, without a second derivative of $A$,
\[
 |[\Delta',h]A|_{(-3)}
 \le2\norm h_\infty C_\Delta(a_0+a_1)|A|_{(-1)}.
\]
After composition with $G_\square h$, the bound is
$2\norm h_\infty^2 C_\Delta(a_0+a_1)C_G$.
This gives no inverse-block-size gain; a bounded constant is enough
for the head summation principle.

An operator bound is not automatically an absolute kernel-row bound.
Here the finite stencil supplies the missing step: testing on one
input bounds each weighted matrix coefficient, and summing costs
at most the number $m_d$ of stencil bonds. For the indicated one-step
locality, the conservative choice $m_d=d5^d$ covers all orientations
in a radius-two cube. Exponential weights add their maximal stencil
ratio. This proves the required coefficientwise reduction, subject
to the same local inverse and retained-weight comparisons. The local
$P_1$ terms, complex averaging estimate, and sum over all heads are
still separate inputs to the global current argument.

\section{Sharp defects for a complex Gaussian}

Let $\Lambda$ be a finite subset of a countable metric space. Each
$i\in\Lambda$ has a real variable $x_i\in\R^D$, with fixed $D$.
Let $C=C^T$ be a complex symmetric covariance on
$\C^{D|\Lambda|}$, with
\begin{equation}\label{eq:cov-assumptions}
 \Rea C\ge cI,\qquad \norm C\le C_*,\qquad c>0,
\end{equation}
where real part denotes the Hermitian part. Gaussian integration is
over the real coordinate space, with the analytic determinant branch
continued from positive real covariances. For fixed measurable
functions $q_i$, assume
\begin{equation}\label{eq:defect}
 |q_i(x_i)|\le\mathbf1_{\{|x_i|>R\}}.
\end{equation}
The sharp cutoff is $q_i=-\mathbf1_{\{|x_i|>R\}}$. Its threshold and
its real integration domain are fixed as the covariance varies.

\begin{lemma}\label{lem:tail}
For every nonempty $S\subset\Lambda$,
\[
 \left|\E_C\prod_{i\in S}q_i\right|
 \le\left[
  (C_*/c)^D 2^{D/2}
  \exp\!\left(-\frac{cR^2}{4C_*^2}\right)
 \right]^{|S|}.
\]
This expectation is holomorphic in $C$ on the accretive cone.
\end{lemma}

\begin{proof}
The marginal covariance $C_{SS}$ inherits both bounds in
\eqref{eq:cov-assumptions}. In $N=D|S|$ coordinates,
\[
 \Rea C_{SS}^{-1}
 =C_{SS}^{-*}(\Rea C_{SS})C_{SS}^{-1}
 \ge cC_*^{-2}I,
 \qquad |\det C_{SS}|\ge c^N.
\]
The absolute Gaussian density is therefore bounded by $(C_*/c)^N$
times the real isotropic Gaussian of variance $C_*^2/c$.
For a $D$-dimensional block of variance $v$,
Markov's inequality applied to $e^{|x|^2/(4v)}$ gives
$\Pr(|x|>R)\le2^{D/2}e^{-R^2/(4v)}$.
That comparison Gaussian has independent blocks, proving the bound
without independence of the original variables.
Locally uniform Gaussian domination proves holomorphy under the
integral. No derivative of $q_i$ is taken.
\end{proof}

\begin{lemma}[A useful precision formulation]\label{lem:precision}
If $A=A^T$, $\Rea A\ge aI$, $\norm A\le B$, and $C=A^{-1}$, then
$\Rea C\ge aB^{-2}I$ and $\norm C\le a^{-1}$. Moreover,
\[
 \left|\E_{A^{-1}}\prod_{i\in S}q_i\right|
 \le\left[(B/a)^D2^{D/2}e^{-aR^2/4}\right]^{|S|}.
\]
\end{lemma}

\begin{proof}
Let $P=S$, $Q=\Lambda\setminus S$, and let
$S_A=A_{PP}-A_{PQ}A_{QQ}^{-1}A_{QP}=(C_{PP})^{-1}$.
For $v=(u,-A_{QQ}^{-1}A_{QP}u)$, one has $Av=(S_Au,0)$.
Hence $\Rea\langle u,S_Au\rangle\ge a\norm u^2$.
The lower bound on $\Rea C_{PP}$ also gives $\norm{S_A}\le B^2/a$.
The absolute marginal density is bounded by $(B/a)^{D|S|}$ times
the real Gaussian of precision $aI$. Apply the same block tail bound.
\end{proof}

\section{From marked defects to connected polymer activities}

For $i\ne j$ put $a_{ij}=\norm{C_{ij}}$ and $a_{ii}=0$.
Fix $\kappa\ge0$ and assume
\begin{equation}\label{eq:row-sum}
 A_\kappa=\sup_i\sum_{j\ne i}a_{ij}e^{\kappa d(i,j)}<\infty.
\end{equation}
Let $\ell(S)$ denote the minimum spanning-tree length on $S$, with
$\ell(\{i\})=0$. Define safe constants
\begin{equation}\label{eq:constants}
 c'=c/2,\quad C'=C_*+c/2,\quad \eta_0=c/4,\quad
 K=\frac{2e^2}{\eta_0}=\frac{8e^2}{c},\quad
 t_R=(C'/c')^D2^{D/2}e^{-c'R^2/(4C'^2)}.
\end{equation}
The symbol $\eta_0$ here is a covariance perturbation radius, unrelated
to lattice spacing in Sections 1--2.

\begin{theorem}\label{thm:polymer}
There are activities $\zeta(S)$ for all nonempty $S\subset\Lambda$
such that
\begin{equation}\label{eq:polymer-gas}
 Z_\Lambda=\E_C\prod_{i\in\Lambda}(1+q_i)
 =\sum_{\{S_1,\ldots,S_m\}\text{ pairwise disjoint}}
   \prod_{j=1}^m\zeta(S_j),
\end{equation}
where the empty family contributes one. For $|S|=n\ge2$,
\begin{equation}\label{eq:tree-bound}
 |\zeta(S)|\le t_R^nK^{n-1}
 \sum_{T\text{ tree on }S}
   \prod_{i\in S}d_i(T)!\prod_{\{i,j\}\in T}a_{ij},
\end{equation}
and $|\zeta(\{i\})|\le t_R$. For every $\alpha>0$, set
\[
 r=\frac{27}{4}e^\alpha t_RKA_\kappa.
\]
If $r<1$, then
\begin{equation}\label{eq:anchored-activities}
 \sup_x\sum_{S\ni x}|\zeta(S)|
     e^{\alpha|S|+\kappa\ell(S)}
 \le\frac{e^\alpha t_R}{1-r}.
\end{equation}
All activities are holomorphic in any background on which $C$ is
holomorphic and the hypotheses hold uniformly.
\end{theorem}

\begin{proof}
For a marked set $S$, apply the BKAR forest formula to
$F_S(s)=\E_{C(s)}\prod_{i\in S}q_i$, where
$C(s)_{ij}=s_{ij}C_{ij}$ for $i\ne j$ and the diagonal is unchanged.
The formula is needed only on the convex hull of partition masks
\cite{bkar,abdriv}; pinching by each partition preserves
$\Rea C\ge cI$ and contracts operator norm. Thus every interpolated
covariance has the same accretivity and norm bounds. No assumption
about arbitrary masks in the full parameter cube is needed.

For a tree $T$, its interpolated mask has entry equal to the minimum
edge parameter along the unique tree path. Write this mask as $s^T(u)$.
The tree activity is
\begin{equation}\label{eq:tree-activity}
 \zeta(S)=\sum_{T\text{ tree on }S}
 \int_{[0,1]^T}\partial_TF_S(s^T(u))\,du.
\end{equation}
For a singleton it is the one-block expectation. Forest components
factor because their interpolated cross covariances vanish. Therefore
$\E_C\prod_{i\in S}q_i$ is the sum over partitions of $S$ of the
products of these activities. Expanding $\prod(1+q_i)$ proves
\eqref{eq:polymer-gas}.

We now bound the derivatives in \eqref{eq:tree-activity} without
differentiating any indicator. Let $d_i=d_i(T)$ and, for each nonzero
edge block, use the complex parameter radius
\[
 \rho_{ij}=\frac{\eta_0}{(d_i+d_j)a_{ij}}.
\]
A simultaneous perturbation in these edge parameters changes each
block row sum by at most
$\eta_0\sum_{j:\{i,j\}\in T}(d_i+d_j)^{-1}\le\eta_0$.
The same holds for columns, so the operator norm of the perturbation
is at most $\eta_0$. The complete polydisc remains within the safer
bounds $\Rea C\ge c'I$, $\norm C\le C'$.
Lemma~\ref{lem:tail} gives $|F_S|\le t_R^n$ there. Cauchy's formula
in the $n-1$ distinct tree variables yields
\[
 |\partial_TF_S|
 \le t_R^n\eta_0^{-(n-1)}
       \prod_{\{i,j\}\in T}(d_i+d_j)a_{ij}.
\]
An edge with $a_{ij}=0$ has zero derivative and is simply omitted.
Since $d_i+d_j\le2d_id_j$, and
$d_i^{d_i}\le e^{d_i}d_i!$, the last product is bounded by
$(2e^2)^{n-1}\prod_i d_i!\prod_{\{i,j\}\in T}a_{ij}$.
This proves \eqref{eq:tree-bound}.

For completeness, the exact degree-weighted tree sum is
\begin{equation}\label{eq:prufer}
 \sum_{T\text{ on }[n]}\prod_i d_i!
 =(n-2)!\binom{3n-3}{n-2}.
\end{equation}
Indeed, the Pr\"ufer multiplicities $k_i=d_i-1$ sum to $n-2$.
For fixed multiplicities, multiplication of the number of trees by
$\prod_i d_i!$ gives $(n-2)!\prod_i(k_i+1)$.
Sum using the coefficient of $z^{n-2}$ in $(1-z)^{-2n}$.

Fix the first vertex at $x$. Replace unordered marked sets by labelled
positions with factor $1/(n-1)!$, drop distinctness, and sum leaf
positions successively. The weighted kernel sum at each removal is
at most $A_\kappa$. Since $\ell(S)\le\sum_{e\in T}d(e)$,
\[
 \sum_{S\ni x,\ |S|=n}|\zeta(S)|e^{\kappa\ell(S)}
 \le t_R^n(KA_\kappa)^{n-1}
       \frac{1}{n-1}\binom{3n-3}{n-2}.
\]
For $m=n-1$, the binomial is at most
$\binom{3m}{m}\le(27/4)^m$; the latter follows by bounding one
term in $(1+1/2)^{3m}$. After inserting $e^{\alpha n}$ and summing,
the bound is $e^\alpha t_R[1+\sum_{m\ge1}r^m/m]$, which is at most
\eqref{eq:anchored-activities}.
\end{proof}

\section{A zero-free logarithm with an anchored bound}

\begin{corollary}\label{cor:log}
Under Theorem~\ref{thm:polymer}, suppose
\begin{equation}\label{eq:smallness}
 r\le\tfrac12,\qquad 2e^\alpha t_R\le\alpha.
\end{equation}
Then $Z_\Lambda\ne0$, with a logarithm continued from zero defect
activity, and
\[
 \log Z_\Lambda=\sum_{\varnothing\ne X\subset\Lambda}e_\Lambda(X),
 \qquad
 \sup_x\sum_{X\ni x}e^{\kappa\ell(X)}|e_\Lambda(X)|
 \le2e^\alpha t_R.
\]
The expansion is absolutely convergent and holomorphic uniformly on
any background domain with the same constants. The subscript $\Lambda$
is retained: consistency between different finite-volume covariances
has not been assumed.
\end{corollary}

\begin{proof}
Put $\rho(S)=|\zeta(S)|e^{\kappa\ell(S)}$. Polymers are incompatible
exactly when they overlap, including self-incompatibility. Equation
\eqref{eq:anchored-activities} gives
\[
 \sum_{T:T\cap S\ne\varnothing}\rho(T)e^{\alpha|T|}
 \le\alpha|S|.
\]
The Koteck\'y--Preiss criterion, in its pinned-series formulation
\cite{kp,fp}, therefore bounds the absolute cluster sum pinned at $S$
by $e^{\alpha|S|}$.

Define $e_\Lambda(X)$ by collecting terms of the usual logarithmic
cluster series whose union is $X$; repeated polymers in that series
are allowed. A connected overlap graph makes the union of spanning
trees inside its polymers connected, so
$\ell(X)\le\sum_j\ell(S_j)$. For a cluster containing $x$,
bound its indicator by $\sum_j\mathbf1_{\{x\in S_j\}}$ and root
at that polymer. The factorial in the cluster series then gives
\[
 \sum_{X\ni x}e^{\kappa\ell(X)}|e_\Lambda(X)|
 \le\sum_{S\ni x}\rho(S)e^{\alpha|S|}
 \le2e^\alpha t_R.
\]
Absolute convergence proves the identity with the finite partition
function, nonvanishing, and the stated holomorphy.
\end{proof}

\begin{remark}[Background derivatives]
If $C$ is Fr\'echet holomorphic and these hypotheses hold uniformly
on an open background ball of radius $r_0$, the
family $\{e_\Lambda(X)\}$ is holomorphic with values in its anchored
weighted sequence space. Cauchy's formula, using radii
$r_0/(2\norm h)$ and $r_0/(2\norm k)$, bounds the mixed second
derivative at the center in that norm by
$8e^\alpha t_R\norm h\norm k/r_0^2$.
The boundary radii are approached from below. This statement concerns
each finite $\Lambda$ with a common bound; no assertion of an
infinite-volume background functional is implicit.
This preserves smallness but gives no additional spatial localization
in the supports of $h$ and $k$.
\end{remark}

\subsection{Logarithmic real cutoffs}

For fixed constants and $R(g)=A(\log g^{-2})^p$, $p>1/2$, one has
$t_{R(g)}=C_0\exp[-c_0(\log g^{-2})^{2p}]$. Consequently, for
every fixed $N>0$ and real $q$,
\[
 t_{R(g)}=o\!\left(g^N(\log g^{-2})^q\right)
 \quad(g\downarrow0).
\]
To check this, divide logarithms by $(\log g^{-2})^{2p}$: the terms
of order $\log g^{-2}$ and $\log\log g^{-2}$ vanish. Thus sufficiently
small $g$ enforces \eqref{eq:smallness}. This conclusion is uniform
in scale only if $D,c,C_*,A_\kappa,A,p$ are themselves uniform.
It is a numerical consequence, not a proof that the actual background
domains have those constants.

\section{A direct quadratic bound without background-support activities}

For the scalar running-coupling correction, it can suffice to estimate
the total Gaussian cutoff Hessian. This is a narrower task than
constructing local effective actions for the complete RG induction.

Assume a metric lattice with volume-independent sums
$S_\gamma=\sup_x\sum_i e^{-\gamma d(x,i)}<\infty$ for
$\gamma>0$. In addition to \eqref{eq:cov-assumptions}, assume
$\norm{C_{ij}}\le K_0e^{-\mu d(i,j)}$. Put
\[
 \nu=\mu/2,\qquad 0<\kappa<\nu,\qquad
 \rho=\frac{c}{8S_\nu}.
\]
Two symmetric covariance perturbations, each with entries bounded by
$e^{-\nu d(i,j)}$, have combined operator norm at most $c/4$ when
their two complex coefficients have modulus at most $\rho$.
Thus the complete polydisc has the common safe bounds
\[
 \Rea C\ge c/2,
 \quad\norm C\le C_*+c/2,
 \quad\norm{C_{ij}}\le (K_0+2\rho)e^{-\nu d(i,j)}.
\]
Apply Theorem~\ref{thm:polymer} with these constants. Its internal
Cauchy margin gives the tail factor
\[
 t_R^*=2^{D/2}
 \left(\frac{C_*+3c/4}{c/4}\right)^D
 \exp\!\left[-\frac{(c/4)R^2}{4(C_*+3c/4)^2}\right].
\]
Choose $R$ large enough that the two smallness conditions
\eqref{eq:smallness} hold with these safe constants and $\alpha=1$.
The tree weights in \eqref{eq:tree-bound} can now all be replaced by
the fixed envelope $(K_0+2\rho)e^{-\nu d(i,j)}$. Their positive
logarithmic cluster series defines one common majorant $M(X)$ such
that, throughout every such polydisc,
\begin{equation}\label{eq:common-majorant}
 |e_\Lambda(X;C)|\le M(X),\qquad
 \sup_z\sum_{X\ni z}e^{\kappa\ell(X)}M(X)
 \le M_0:=2e\,t_R^*.
\end{equation}
This is a majorant for the sum of suprema, rather than a separate
estimate whose supremum might depend on $X$.
The union coefficient $e_\Lambda(X;C)$ reads only $C_{XX}$.

\begin{theorem}\label{thm:direct-hessian}
Let $C(s,t)$ be a holomorphic symmetric covariance with value $C$
at the origin. Write $E_x=\partial_sC$, $E_y=\partial_tC$, and
$E_{xy}=\partial_s\partial_tC$ there. Suppose, for each nonempty
$X$ and $i,j\in X$,
\begin{align*}
 \norm{(E_x)_{ij}}&\le B_1e^{-\mu\operatorname{dist}(x,X)}
                            e^{-\nu d(i,j)},\\
 \norm{(E_y)_{ij}}&\le B_1e^{-\mu\operatorname{dist}(y,X)}
                            e^{-\nu d(i,j)},\\
 \norm{(E_{xy})_{ij}}&\le B_2
 e^{-\nu[d(x,y)+\operatorname{dist}(x,X)+\operatorname{dist}(y,X)]}
 e^{-\nu d(i,j)}.
\end{align*}
Then the Gaussian cutoff logarithm $F_\Lambda(C)=\log Z_\Lambda(C)$
satisfies
\begin{equation}\label{eq:direct-hessian}
 |\partial_s\partial_tF_\Lambda(C(s,t))|_{s=t=0}
 \le M_0\left[
 \frac{B_1^2}{\rho^2}S_{\mu-\kappa}
 +\frac{B_2}{\rho}S_\nu\right]e^{-\kappa d(x,y)}.
\end{equation}
\end{theorem}

\begin{proof}
For directions on $X$ with entry envelopes
$b_Ee^{-\nu d(i,j)}$ and $b_Fe^{-\nu d(i,j)}$, extend by zero
outside $X\times X$. Cauchy's formula with radii $\rho/b_E$ and
$\rho/b_F$, and the common majorant, gives
\[
 |De_\Lambda(X)[E]|\le M(X)b_E/\rho,
 \qquad
 |D^2e_\Lambda(X)[E,F]|\le M(X)b_Eb_F/\rho^2.
\]
Each variable is differentiated once, so no factor two is missing.
With $\delta_x=\operatorname{dist}(x,X)$ and similarly for $y$,
the chain rule bounds the summand by
\[
 M(X)\left[
 \frac{B_1^2}{\rho^2}e^{-\mu(\delta_x+\delta_y)}
 +\frac{B_2}{\rho}e^{-\nu[d(x,y)+\delta_x+\delta_y]}
 \right].
\]
The elementary inequality
$d(x,y)\le\delta_x+\ell(X)+\delta_y$ implies
\[
 e^{-\mu(\delta_x+\delta_y)}
 \le e^{-\kappa d(x,y)}e^{\kappa\ell(X)}
       e^{-(\mu-\kappa)\delta_x}.
\]
For $q>0$, $e^{-q\delta_x}\le\sum_{z\in X}e^{-qd(x,z)}$.
Equation \eqref{eq:common-majorant} therefore makes the first sum
at most $M_0S_{\mu-\kappa}e^{-\kappa d(x,y)}$ times its coefficient.
For the second, use $\nu>\kappa$, drop the $y$-distance factor,
and sum with $S_\nu$. This proves \eqref{eq:direct-hessian}.
\end{proof}

\subsection{Local precision variations provide the response bounds}

Let $A=C^{-1}$ and $A(s,t)=A+sH_x+tH_y$, where each symmetric
$H_x$ is supported on $B_{r_s}(x)\times B_{r_s}(x)$, has norm at
most $H$, and the ball has at most $n_s$ sites. Inverse differentiation
gives
\[
 E_x=-CH_xC,\qquad
 E_{xy}=CH_xCH_yC+CH_yCH_xC.
\]
Expanding in blocks yields
\begin{align*}
 \norm{(E_x)_{ij}}
 &\le b_1e^{-\mu[d(i,x)+d(j,x)]},
 &b_1&=n_s^2HK_0^2e^{2\mu r_s},\\
 \norm{(CH_xCH_yC)_{ij}}
 &\le b_2e^{-\mu[d(i,x)+d(x,y)+d(y,j)]},
 &b_2&=n_s^4H^2K_0^3e^{4\mu r_s}.
\end{align*}
The first distance sum is at least both $d(i,j)$ and
$2\operatorname{dist}(x,X)$. The second is at least both $d(i,j)$
and $d(x,y)+\delta_x+\delta_y$. Averaging these two lower bounds
in each case verifies Theorem~\ref{thm:direct-hessian} with
$B_1=b_1$ and $B_2=2b_2$. Source supports may overlap, and $x=y$
is allowed.

For a nonlinear precision the exact second response has the additional
term $-C(\partial_x\partial_yA)C$. It must be estimated too.
For example, if that precision derivative vanishes for $d(x,y)>r_2$
and otherwise has uniformly bounded norm and support in a fixed ball
about $x$, the same first-variation calculation supplies an extra
constant times $t_R^*e^{\kappa r_2}e^{-\kappa d(x,y)}$.
No locality assumption on first variations alone controls this term.

If a Hessian kernel bounded by $C_Ht_R^*e^{-\kappa d(z,w)}$ is
contracted with responses
$|h_x(z)|\le H_he^{-m d(x,z)}$, then, for
$\lambda=\frac12\min(\kappa,m)$, the contraction is bounded by
\[
 C_Ht_R^*H_h^2S_{m-\lambda}S_{\kappa-\lambda}
 e^{-\lambda d(x,y)}.
\]
This follows by the triangle inequality through $x,z,w,y$ and two
geometric sums. Differentiating a nonlinear response chart instead
adds its own second-variation term; the displayed contraction alone
does not include it.

\subsection{Thermodynamic limit and the scalar coefficient}

Suppose, additionally, that finite-volume covariances and their first
and second source responses converge entrywise on every fixed set,
with the preceding bounds uniform. Fixed-set Hessians then converge:
the union coefficient has the finite inclusion--exclusion formula
\[
 e_\Lambda(X;C)=\sum_{Y\subset X}(-1)^{|X\setminus Y|}
       \log Z_Y(C_{YY}),
\]
where the empty-set logarithm is zero. The uniform zero-free domain
makes these finite-dimensional functions and their derivatives
continuous under the indicated convergence.
For sets $X$ not contained in $B_L(x)$,
$\operatorname{dist}(x,X)+\ell(X)\ge L$.
Spending, for example, $\kappa/2$ of the available tree and source
decay gives an omitted Hessian sum bounded by
$C t_R^*e^{-\kappa L/2}$, uniformly in volume. Thus the total Hessian
converges too. The argument applies to growing tori when fixed balls
embed with their ordinary lattice metric and all counting bounds are
uniform. Entrywise convergence is an additional hypothesis, not a
consequence of the uniform estimate.

At finite volume, any differentiable analytic pointing identity
\[
 F_{{\rm cut},\Lambda}(\xi)-F_{{\rm cut},\Lambda}(0)
 =\sum_{X,z}E_{{\rm cut},\Lambda}(X;\xi,z)
\]
identifies the total Hessian with the sum of its pointed Hessians.
Therefore this direct route can bound the pure Gaussian cutoff part
of the deposited quadratic coefficient after its normalized response
and infinite-volume kernel have been identified. On $\mathbb Z^4$,
\[
 \sum_y d(x,y)^2|D_xD_yF_{\rm cut}|
 \le C t_R^*\sum_y d(x,y)^2e^{-\lambda d(x,y)}<\infty.
\]
In particular the cutoff contribution to the coupling increment is
$O(t_R^*)$ under those identifications.

The exact decomposition is
\[
 \log\E_C[\chi_R e^V]
 =\log\E_C[\chi_R]
 +\log\frac{\E_C[\chi_R e^V]}{\E_C[\chi_R]}.
\]
This section bounds the first term. The interaction ratio and the
separate Gaussian normalization remain to be controlled. It also
does not replace background-local effective actions in the full RG
induction. It shows that their complete construction is unnecessary
for this particular quadratic estimate once the direct source
responses and kernel identification are proved.

\section{The remaining Yang--Mills interface}

\subsection{Matching the actual Gaussian variables}

The independent variables are identifiable in the primary construction.
CMP109, pp.\ 266--268, (2.9)--(2.13), linearizes the averaging
constraint, rescales by the running coupling, and removes a pivot bond
per next-level bond. The remaining bond vector $B$ is unconstrained;
its precision is $\mathcal C^*\Delta^{(k)}\mathcal C$, where
$\mathcal C$ denotes the extension map, not a covariance. Each retained
bond has $D=\dim_\R\mathfrak g$ coordinates. The paper explicitly
changes the earlier delta gauge to exponential gauge
\cite{balaban109}.

CMP99b, p.\ 428 after (3.158), gives a scale-independent positive lower
bound for the axial reduced precision; Theorem 3.15, p.\ 432, gives
exponential covariance decay and walks in its stated regular-background
setting \cite{balaban99}. These are relevant inputs, but the axial and
exponential gauge matrices have different coordinates. Moreover,
the primary current is tied to the background, whereas the later
analytic carrier permits an independent current.

CMP116, pp.\ 12--14, (2.1)--(2.10), includes the cutoff-only sector
when the interaction sets are empty and the bad-bond set is nonempty.
Its localization uses conditioning, exterior whitening, and walks.
The complex estimates at pp.\ 15--17 are structured domination
estimates, not an explicitly stated accretive-covariance theorem on
every arbitrary complex walk polydisc \cite{balaban116}.
Our forest interpolation is a separate, accretivity-preserving device.

The deposit already offers conditional operator constructions here.
Raw I, lines 5910--6077 and 7171--7298, derives accretivity and
localized covariance bounds from named form-walk and positivity
inputs. Lines 5500--5517 distinguish primary tied-current estimates
from the free-current hypotheses. Thus the missing task is to verify
the hypotheses and coordinate identification, not to pretend that
no proposed operator construction exists.

\begin{proposition}[A sufficient uniform perturbation criterion]
For $\mu\ge0$, let $\norm{\cdot}_\mu$ be the maximum of the exponentially weighted
block row and column sums. Suppose the real reference precision $A_0$
is symmetric, $A_0\ge2aI$, $\norm{A_0}\le B_0$, and
$\norm{A_0^{-1}}_\mu\le K_0$. Let the symmetric holomorphic
perturbation $V(\xi)$ on a fixed background/current ball satisfy
\[
 \norm{V(\xi)}_\mu\le\delta,\qquad
 \delta\le a,\qquad K_0\delta<1.
\]
Then $A=A_0+V$ has $\Rea A\ge aI$, $\norm A\le B_0+\delta$, and
\[
 \norm{A^{-1}}_\mu\le\frac{K_0}{1-K_0\delta}.
\]
In particular the covariance hypotheses of
Theorem~\ref{thm:polymer} hold with constants independent of scale
whenever the displayed inputs are independent of scale.
\end{proposition}

\begin{proof}
The weighted row/column norm is a matrix algebra norm by the triangle
inequality for distance, and dominates operator norm by the Schur
test. Accretivity and the upper bound follow immediately. Sum
$A^{-1}=(I+A_0^{-1}V)^{-1}A_0^{-1}$ in that algebra.
Lemma~\ref{lem:precision} supplies covariance accretivity.
\end{proof}

This criterion reduces the analytic extension to a concrete weighted
operator perturbation estimate on the \emph{actual} independent
coordinates. It does not prove that estimate or replace background
locality.
Notebook 06 now proves both real flat precision bounds directly for
the explicit full-block Maxwell quadratic model and gives the exact
minimizer/coordinate interface to the physical Gaussian. Notebook 07
adds exponential flat kernel decay and an $O(M^{-2})$ refinement limit
in weighted kernel norm, uniformly over compatible periodic volumes.
This verifies the homogeneous flat reference inputs of the criterion.
The weighted bound on the actual complex perturbation is still needed;
neither spectral bounds nor spatial decay alone provides it.
Notebook 08 additionally proves local fine response legs for the
normalized homogeneous flat source minimizer: the potential has
uniform physical $C^{1,\alpha}$ bounds, its curl has $C^\alpha$
bounds, and its componentwise Laplacian has a uniform supremum
bound, all with exponential coarse-cell decay. The corrected
averaging map, scalar gauge projection and normalization are
checked explicitly. Notebook 09 now also proves locality of the full
flat fine inverse and a common complex neighborhood for the deposited
auxiliary current $G_\pi^J$. It supplies a quadratic zero-current
precision remainder and all-order localized current derivatives,
uniformly in refinement and compatible periodic volume. Notebook 11
also controls the separate nonlinear-constraint term by its corrected
fixed-$L$ quadratic composition, and proves a full flat $G_1^J$
current neighborhood in the physical-cell algebra for $M=L^j$.
Its passage to the actual retained Gaussian precision and nonflat
background domains remain to be matched. Notebook 10 controls small curvature-metric
perturbations at fixed $Q$, which is a further quadratic input and
does not identify an entire nonflat Hessian.
Notebook 12 additionally proves the full flat current bounds in
the homogeneous pointwise $N,N_\beta$ norms for bounded fine
forcing, with no derivative of the current. This is a further
matched source input. Its additional heat-kernel argument handles
the one-form right-divergence data, with equal exponents for every
fixed H\"older exponent strictly between zero and one,
and a separate series handles the constrained-current inverse.
The scalar member has a separate proof under an explicit uniform
positive-penalty hypothesis. Strict dependence on background
supports and the retained-coordinate comparison remain distinct
obligations. Notebook 13 further straightens the complete nonlinear
fixed-reference constraint on a refinement-uniform analytic ball;
Notebook 14 computes the actual moving-reference derivative
correction, controls the ordinary reference averages, and proves
weighted energy-inverse stability for a Wilson operator on a small
global physical $C^1$ ball. Its bare action, physical pairing and
left-exponential Hessian chart match the deposited conventions in
dimension four. Original-paper verification and the precise
graded-curl and effective-Hessian identifications remain separate.
Notebook 15 now constructs a nonlinear action minimizer in a
homogeneous small-global-field class, with uniform physical
$C^{1,\alpha}$ control and analyticity. A gauge flow and an exact
ordinary-average transfer remove the fixed-slice and corrected-map
restrictions for its local action comparison. Nonflat source
pointwise/current domains, mixed retained geometry, and the actual
fluctuation integration remain separate interfaces.
Notebook 16 further proves weighted responses for the corrected
and explicitly chosen ordinary minimizers. Its classical action
has derivative equal to the coarse multiplier and a nonlinear
Hessian that is a small exponentially local perturbation of the
flat Maxwell precision. These estimates concern a classical
coarse Hessian, with the ordinary-coordinate convention explicit;
they do not identify the precision after quantum integration or
give strict background-support activities.

Notebook 17 now proves joint correlated bad-block bounds, including
a complex local-event estimate without a total-volume loss. It
compares the first-step reduced Wilson Hessian with the flat energy
and identifies the homogeneous flat retained Gaussian at all earlier
averaging depths using the nonlinear minimum of Notebook 15.
Notebook 18 closes the abstract overlapping-interaction expansion
under an explicit weighted coefficient norm; with a small local
precision decomposition, its activities also have strict background
support. Notebook 19 now identifies the actual physical homogeneous
pivot Gaussian and bounds its local measure and gauge factors.
Notebook 20 matches all local interaction and precision hypotheses
for the bare first-step Wilson small-field integral, including
its normalization terms. Notebook 21 proves the support-valued
classical action and precision bounds uniformly in earlier
refinement, and closes one integration of that classical reference
action. Notebook 22 constructs the strict stationary-background
support expansion and transports those activities to the next
coarse variables with explicit radius and support-weight losses.
Notebook 23 proves local gauge covariance of the full correction
and its quadratic Ward identity. Notebook 24 proves a conditional
inherited fluctuation-interaction estimate, and Notebook 26
meets its support-weight test using fixed auxiliary cells.
Notebook 25 supplies a smaller classical gauge representative
and an analytic composition-germ continuation. Stable quantum
domains and amplitudes, including the marginal coupling flow,
remain to be proved for the induction.

\subsection{Marked-set decay is not background locality}

The preceding activities involve the marked locations in $S$, but
$C_{SS}(W)$ may depend on $W$ arbitrarily far away. Consider, for
example,
\[
 C(W)=(1+\varepsilon W_{x_*})I,\qquad
 |\varepsilon W_{x_*}|<\tfrac14.
\]
It satisfies uniform accretivity and norm bounds and has no
off-diagonal covariance at all. Nevertheless, every singleton cutoff
expectation depends on the remote background coordinate $W_{x_*}$.
Therefore Theorem~\ref{thm:polymer} alone cannot supply a family of
activities supported in their claimed background domains.

Notebook 04 now proves a sufficient strict-locality theorem when the
precision is a sufficiently small sum of local perturbations of a fixed,
uniformly accretive, exponentially decaying baseline. It expands the
entire marginal density, including its determinant and inverse, and
pays for roots in the added support. Its product-domain, symmetry,
finite-volume, and actual-operator interfaces remain explicit.

For the actual Yang--Mills Gaussian, the needed further input is a
local walk expansion, or an equivalent
analytic decoupling construction, for the \emph{background dependence}
of the actual Gaussian. It must control the weighted sums of the
domains visited by those walks. A truncation to balls with error
$e^{-cR}$ does not automatically pay a polymer norm that penalizes
their four-dimensional volume. Source-localization and volume
compatibility must be checked in the norm actually used by the RG.

\subsection{What has to be matched to the deposited construction}

The current proof and the Gaussian proof leave a short list of exact
interfaces rather than an unidentified analytic remainder:
\begin{enumerate}
\item Prove a weighted bound on the combined global current
$s_\sigma$, on the same complex domains used by the four-stage
induction. If the earlier substituted representation is used,
estimate \eqref{eq:current-residual} and prove the additional global
operator identifications.
\item Propagate the derived-current clauses at every retained triple.
Weak membership of the basic carrier alone does not establish these.
\item Identify the actual independent Gaussian coordinates and prove
uniform accretivity and decay for their full covariance on the fixed
complex, independent-current domain. The forest interpolants then
inherit these bounds; accretivity of the older complex walk
interpolation is not an extra premise of our Gaussian theorem.
A bound only on physical tied currents is insufficient.
\item For the full effective-action induction, localize the Gaussian
cutoff logarithm in the background variables with an anchored bound
on a fixed normalized complex domain. For the scalar quadratic
coefficient alone, Section 6 gives the alternative of direct
source-response bounds and local convergence. Control the
interaction contribution and separate Gaussian normalization as well.
Notebook 05 supplies the flat gauge/coercivity bridge and a complete
interaction module for disjoint cells. Notebook 18 extends the
integration to overlapping supports and combines it with the local
precision construction. Its weighted coefficient hypotheses must
still be verified for the actual Yang--Mills integrand.
\item Apply the pointed quadratic extraction in Notebook 01 to that
actual local family, with uniform normalized source responses. The
extraction lemma does not establish its own analytic hypotheses.
\end{enumerate}

These tasks concern a necessary renormalization estimate. The mixed
observable hierarchy, all-order continuum bounds, rotation restoration,
identification as Yang--Mills, and the mass gap remain separate
obligations. Notebook 02 supplies a fixed-coupling reflection argument;
it supplies no uniform ultraviolet estimate for this list.

\begin{thebibliography}{9}
\bibitem{bkar} A. Abdesselam, notes on the BKAR formula, Theorem 1
and the convex hull of partition masks.
\href{https://mabdesselam.github.io/BKAR.pdf}{Author-hosted notes}.
\bibitem{abdriv} A. Abdesselam and V. Rivasseau,
\emph{Trees, forests and jungles: a botanical garden for cluster
expansions}, Lecture Notes in Physics 446 (1995), 7--36.
\href{https://arxiv.org/abs/hep-th/9409094}{arXiv:hep-th/9409094}.
\bibitem{kp} R. Koteck\'y and D. Preiss,
\emph{Cluster expansion for abstract polymer models},
Commun. Math. Phys. 103 (1986), 491--498.
\href{https://doi.org/10.1007/BF01211762}{Publisher record}.
\bibitem{fp} R. Fern\'andez and A. Procacci,
\emph{Cluster expansion for abstract polymer models. New bounds from
an old approach}, Commun. Math. Phys. 274 (2007), 123--140,
especially equations (2.7), (2.12), (2.14), and (2.15).
\href{https://webspace.science.uu.nl/~ferna107/papers/ferpro07.pdf}
{Author-hosted paper};
\href{https://arxiv.org/abs/math-ph/0605041}{arXiv version}.
\bibitem{balaban99} T. Ba\l aban,
\emph{Propagators for lattice gauge theories in a background field},
Commun. Math. Phys. 99 (1985), 389--434.
\href{https://doi.org/10.1007/BF01240355}{Publisher record}.
\bibitem{balaban109} T. Ba\l aban,
\emph{Renormalization group approach to lattice gauge field theories I},
Commun. Math. Phys. 109 (1987), 249--301.
\href{https://scispace.com/pdf/renormalization-group-approach-to-lattice-gauge-field-3ygai0vxt1.pdf}
{Copy of the original article}.
\bibitem{balaban116} T. Ba\l aban,
\emph{Renormalization group approach to lattice gauge field theories II.
Cluster expansions}, Commun. Math. Phys. 116 (1988), 1--22.
\href{https://scispace.com/pdf/renormalization-group-approach-to-lattice-gauge-field-4penm4k7hc.pdf}
{Copy of the original article}.
\end{thebibliography}
\end{document}
